\documentclass{article}
\usepackage{macros}
\title{Tarefa 1 - MAT01050}
\author{Leonardo Suarez}
\date{\today}

\setcounter{prblm}{5}

\begin{document}
\maketitle

\begin{prblm}
    Dada \(\displaystyle A = \begin{bmatrix} 0 & 1 & 0 & 0 \\ 0 & 0 & 1 & 0 \\ 0 & 0 & 0 & 1 \\ 2 & 3 & 4 & 5 \end{bmatrix}\), encontre uma matriz de permutação \(P\), uma matriz triangular inferior \(L\) e uma matriz triangular superior \(U\) tais que \(PA = LU\).
\end{prblm}

\begin{proof}[Solução]
    Faz a eliminação gaussiana na matriz \(A\). Troca a primeira e a quarta linha, depois a segunda com a quarta, e a terceira com a quarta. Expressamos essas operações por meio de matrizes elementares \(E_{14}\), \(E_{24}\) e \(E_{34}\), ou seja, temos
    \begin{align}
        P = E_{34}E_{24}E_{14} &= \begin{bmatrix} 1 & 0 & 0 & 0 \\ 0 & 1 & 0 & 0 \\ 0 & 0 & 0 & 1 \\ 0 & 0 & 1 & 0 \end{bmatrix} \begin{bmatrix} 1 & 0 & 0 & 0 \\ 0 & 0 & 0 & 1 \\ 0 & 0 & 1 & 0 \\ 0 & 1 & 0 & 0 \end{bmatrix} \begin{bmatrix} 0 & 0 & 0 & 1 \\ 0 & 1 & 0 & 0 \\ 0 & 0 & 1 & 0 \\ 1 & 0 & 0 & 0 \end{bmatrix} \\
        &= \begin{bmatrix} 1 & 0 & 0 & 0 \\ 0 & 1 & 0 & 0 \\ 0 & 0 & 0 & 1 \\ 0 & 0 & 1 & 0 \end{bmatrix} \begin{bmatrix} 0 & 0 & 0 & 1 \\ 1 & 0 & 0 & 0 \\ 0 & 0 & 1 & 0 \\ 0 & 1 & 0 & 0 \end{bmatrix} \\ 
        &= \begin{bmatrix} 0 & 0 & 0 & 1 \\ 1 & 0 & 0 & 0 \\ 0 & 1 & 0 & 0 \\ 0 & 0 & 1 & 0 \end{bmatrix}.
    \end{align}
    Obtemos então a matriz triangular superior \(U\)
    \begin{equation}
        U = \begin{bmatrix} 2 & 3 & 4 & 5 \\ 0 & 1 & 0 & 0 \\ 0 & 0 & 1 & 0 \\ 0 & 0 & 0 & 1 \end{bmatrix},
    \end{equation}
    e a matriz triangular inferior \(L\), que não possui elementos abaixo da diagonal principal, pois as únicas operações elementares realizadas foram trocas de linha, ou seja,
    \begin{equation}
        L = I = \begin{bmatrix} 1 & 0 & 0 & 0 \\ 0 & 1 & 0 & 0 \\ 0 & 0 & 1 & 0 \\ 0 & 0 & 0 & 1 \end{bmatrix}.
    \end{equation}
\end{proof}

A solução foi verificada no \texttt{Scilab} usando a função \texttt{lu} que retorna as matrizes \(L\), \(U\) e \(P\) a partir da matriz \(A\):
    \begin{verbatim}
        --> A = [[ 0 1 0 0]; [0 0 1 0]; [0 0 0 1]; [2 3 4 5]]

    A = [4x4 double]

    0.   1.   0.   0.
    0.   0.   1.   0.
    0.   0.   0.   1.
    2.   3.   4.   5.

        --> [L,U,P]=lu(A)

    L = [4x4 double]

    1.   0.   0.   0.
    0.   1.   0.   0.
    0.   0.   1.   0.
    0.   0.   0.   1.

    U = [4x4 double]

    2.   3.   4.   5.
    0.   1.   0.   0.
    0.   0.   1.   0.
    0.   0.   0.   1.

    P = [4x4 double]

    0.   0.   0.   1.
    1.   0.   0.   0.
    0.   1.   0.   0.
    0.   0.   1.   0.
    \end{verbatim}

\end{document}